% ======================================================================
%  preamble.tex — Fiches d'automatismes 2G
%  Version unique : \proffalse = élève, \proftrue = professeur
% ======================================================================
\usepackage[T1]{fontenc}
\usepackage[utf8]{inputenc}
\usepackage{mathpazo}
% babel-french est absent de cette installation (TeX Live 2023 allégé) :
% on charge babel en 'english' (motifs de césure anglais, désactivés ci-dessous)
% et la typographie française est assurée à la main (espaces insécables avant ; : ? !,
% guillemets \guillemotleft/\guillemotright, virgule décimale via icomma).
\usepackage[english]{babel}
\usepackage{icomma}
\addto\captionsenglish{\renewcommand{\contentsname}{Sommaire}}
\usepackage[scaled=.92]{helvet}
\usepackage{courier}
\usepackage[a4paper,top=1.7cm,bottom=1.7cm,left=1.9cm,right=1.9cm,headheight=14pt]{geometry}
\usepackage{amsmath,amssymb,mathtools}
\usepackage{xcolor}
\usepackage{enumitem}
\usepackage{needspace}
\usepackage{microtype}
\usepackage{array,booktabs}
\usepackage{tikz}
\usepackage{tkz-tab}
\usepackage[framemethod=default]{mdframed}
\usepackage{fancyhdr}
\usepackage{titlesec}
\usepackage{etoolbox}
\usepackage{stmaryrd}

% --- Palette : marine / acier / ardoise -------------------------------
\definecolor{marine}{HTML}{1F3A5F}
\definecolor{acier}{HTML}{5B7C99}
\definecolor{ardoise}{HTML}{4A5560}
\definecolor{fond}{HTML}{F6F8FA}
\definecolor{rep}{HTML}{8C2F2F}   % réponses (version professeur)

% --- Notations françaises --------------------------------------------
\newcommand{\R}{\mathbb{R}}\newcommand{\N}{\mathbb{N}}\newcommand{\Z}{\mathbb{Z}}
\newcommand{\Q}{\mathbb{Q}}\newcommand{\D}{\mathbb{D}}
\newcommand{\vect}[1]{\overrightarrow{#1}}
\newcommand{\norme}[1]{\left\lVert #1\right\rVert}
\newcommand{\abs}[1]{\left\lvert #1\right\rvert}
\newcommand{\ent}[2]{\llbracket #1\,;#2\rrbracket}      % intervalle d'entiers
\newcommand{\deter}[2]{\det\left(#1,#2\right)}
\newcommand{\Card}{\operatorname{Card}}
\newcommand{\pgcd}{\operatorname{PGCD}}\newcommand{\ppcm}{\operatorname{PPCM}}
\newcommand{\degre}{^{\circ}}
\newcommand{\un}[1]{\,\mathrm{#1}}                      % unité : 3\un{cm}
\newcommand{\sci}[2]{#1\times10^{#2}}
\newcommand{\vc}[2]{\begin{pmatrix}#1\\#2\end{pmatrix}}  % vecteur colonne
\newcommand{\code}[1]{\texttt{#1}}
\setlist{nosep}

% --- Bascule élève / professeur --------------------------------------
\newif\ifprof
\proffalse   % version élève par défaut ; la version professeur s'obtient par
                 % `pdflatex '\def\PROFVERSION{}\documentclass[11pt,a4paper]{article}
% ======================================================================
%  preamble.tex — Fiches d'automatismes 2G
%  Version unique : \proffalse = élève, \proftrue = professeur
% ======================================================================
\usepackage[T1]{fontenc}
\usepackage[utf8]{inputenc}
\usepackage{mathpazo}
% babel-french est absent de cette installation (TeX Live 2023 allégé) :
% on charge babel en 'english' (motifs de césure anglais, désactivés ci-dessous)
% et la typographie française est assurée à la main (espaces insécables avant ; : ? !,
% guillemets \guillemotleft/\guillemotright, virgule décimale via icomma).
\usepackage[english]{babel}
\usepackage{icomma}
\addto\captionsenglish{\renewcommand{\contentsname}{Sommaire}}
\usepackage[scaled=.92]{helvet}
\usepackage{courier}
\usepackage[a4paper,top=1.7cm,bottom=1.7cm,left=1.9cm,right=1.9cm,headheight=14pt]{geometry}
\usepackage{amsmath,amssymb,mathtools}
\usepackage{xcolor}
\usepackage{enumitem}
\usepackage{needspace}
\usepackage{microtype}
\usepackage{array,booktabs}
\usepackage{tikz}
\usepackage{tkz-tab}
\usepackage[framemethod=default]{mdframed}
\usepackage{fancyhdr}
\usepackage{titlesec}
\usepackage{etoolbox}
\usepackage{stmaryrd}

% --- Palette : marine / acier / ardoise -------------------------------
\definecolor{marine}{HTML}{1F3A5F}
\definecolor{acier}{HTML}{5B7C99}
\definecolor{ardoise}{HTML}{4A5560}
\definecolor{fond}{HTML}{F6F8FA}
\definecolor{rep}{HTML}{8C2F2F}   % réponses (version professeur)

% --- Notations françaises --------------------------------------------
\newcommand{\R}{\mathbb{R}}\newcommand{\N}{\mathbb{N}}\newcommand{\Z}{\mathbb{Z}}
\newcommand{\Q}{\mathbb{Q}}\newcommand{\D}{\mathbb{D}}
\newcommand{\vect}[1]{\overrightarrow{#1}}
\newcommand{\norme}[1]{\left\lVert #1\right\rVert}
\newcommand{\abs}[1]{\left\lvert #1\right\rvert}
\newcommand{\ent}[2]{\llbracket #1\,;#2\rrbracket}      % intervalle d'entiers
\newcommand{\deter}[2]{\det\left(#1,#2\right)}
\newcommand{\Card}{\operatorname{Card}}
\newcommand{\pgcd}{\operatorname{PGCD}}\newcommand{\ppcm}{\operatorname{PPCM}}
\newcommand{\degre}{^{\circ}}
\newcommand{\un}[1]{\,\mathrm{#1}}                      % unité : 3\un{cm}
\newcommand{\sci}[2]{#1\times10^{#2}}
\newcommand{\vc}[2]{\begin{pmatrix}#1\\#2\end{pmatrix}}  % vecteur colonne
\newcommand{\code}[1]{\texttt{#1}}
\setlist{nosep}

% --- Bascule élève / professeur --------------------------------------
\newif\ifprof
\proffalse   % version élève par défaut ; la version professeur s'obtient par
                 % `pdflatex '\def\PROFVERSION{}\documentclass[11pt,a4paper]{article}
\input{preamble}   % bascule \proffalse / \proftrue : voir preamble.tex
\begin{document}
\begin{center}
  {\Large\sffamily\bfseries\color{marine} Rituels de début de cours — Mathématiques 2G}\\[3pt]
  {\sffamily\color{ardoise} Année scolaire 2026--2027 — \ifprof Version professeur\else Version élève\fi}
\end{center}
\vspace{2pt}
\begin{mdframed}[style=fichestyle]
\ifprof
  Une fiche par séance de 2G (88 fiches, du 8 octobre 2026 au 1\textsuperscript{er} juillet 2027), alignée sur le calendrier ;
  les questions 4 et 5 tiennent compte du test de positionnement du 29 septembre 2026 (voir le guide et l'annexe A).
\else
  Chaque fiche comprend cinq questions à traiter en sept minutes, sans calculatrice sauf mention.
  Les questions 4 et 5 reviennent sur des notions déjà rencontrées : elles te permettent de te remettre en mémoire ce que tu as déjà appris.
  Indique toujours ta méthode, même pour un calcul mental.
\fi
\end{mdframed}
\tableofcontents
\bigskip
\foreach\i in {01,02,03,04,05,06,07,08,09,10,11,12,13,14}{\input{sections/ch\i}}
\ifprof
  \clearpage
  \input{sections/guide}
  \clearpage
  \input{sections/annexe_diag}
\fi
\end{document}
'` ou en remplaçant \proffalse par \proftrue
\ifdefined\PROFVERSION\proftrue\fi

% --- Compteurs (automatiques) ----------------------------------------
\newcounter{fiche}
\newcounter{qfiche}[fiche]

% --- Étiquette de diagnostic (version professeur) --------------------
\newcommand{\tagdiag}[1]{\textcolor{acier}{\footnotesize\textsf{[#1]}}}

% --- Conteneur d'une fiche ------------------------------------------
\mdfdefinestyle{fichestyle}{%
  linewidth=2pt,topline=false,bottomline=false,rightline=false,
  leftline=true,linecolor=marine,innerleftmargin=9pt,innerrightmargin=2pt,
  innertopmargin=4pt,innerbottommargin=5pt,skipabove=7pt,skipbelow=5pt,
  backgroundcolor=white}
\mdfsetup{style=fichestyle}
\makeatletter
\newcommand{\fichehead}[2]{\gdef\fh@date{#1}\gdef\fh@titre{#2}}
\newenvironment{fiche}{%
  \stepcounter{fiche}\Needspace{9\baselineskip}\begin{mdframed}[style=fichestyle]%
  \noindent{\color{marine}\textbf{\textsf{Fiche~\thefiche}}}\hspace{.6em}%
  {\color{ardoise}\small\textsf{\fh@date}}\hspace{.6em}{\itshape\small\fh@titre}\par\vspace{3pt}%
  \begin{enumerate}[label=\textbf{\arabic*.},leftmargin=1.5em,itemsep=3pt,parsep=0pt]}{%
  \end{enumerate}\end{mdframed}}
\makeatother
\newcommand{\role}[1]{{\color{acier}\scriptsize\textsf{#1}}}
\newcommand{\qitem}[5]{%
  \item \role{#1}\ \ifprof\tagdiag{#2}\fi\ #3%
  \ifprof
    \par\smallskip\noindent{\color{rep}\small\textbf{Rép.}\ #4}%
    \if\relax\detokenize{#5}\relax\else
      \par\noindent{\color{ardoise}\footnotesize\textit{#5}}\fi
  \fi}
\newcommand{\choix}[4]{\par\smallskip\noindent\begin{tabular}{@{}*{4}{p{.22\linewidth}}@{}}
  \textbf{A.}~#1 & \textbf{B.}~#2 & \textbf{C.}~#3 & \textbf{D.}~#4\end{tabular}}

% --- Titres, en-têtes ------------------------------------------------
\titleformat{\section}{\Large\sffamily\bfseries\color{marine}}{}{0pt}{}[\vspace{-2pt}{\color{acier}\titlerule}]
\titlespacing*{\section}{0pt}{12pt}{6pt}
\pagestyle{fancy}\fancyhf{}
\renewcommand{\headrulewidth}{.3pt}
\fancyhead[L]{\small\sffamily Rituels de début de cours — 2G}
\fancyhead[R]{\small\sffamily\ifprof Version professeur\else Version élève\fi}
\fancyfoot[C]{\small\thepage}
\setlength{\parindent}{0pt}\raggedright\hyphenpenalty=10000\exhyphenpenalty=10000
\usepackage[hidelinks]{hyperref}  % toujours en dernier (ou presque)
   % bascule \proffalse / \proftrue : voir preamble.tex
\begin{document}
\begin{center}
  {\Large\sffamily\bfseries\color{marine} Rituels de début de cours — Mathématiques 2G}\\[3pt]
  {\sffamily\color{ardoise} Année scolaire 2026--2027 — \ifprof Version professeur\else Version élève\fi}
\end{center}
\vspace{2pt}
\begin{mdframed}[style=fichestyle]
\ifprof
  Une fiche par séance de 2G (88 fiches, du 8 octobre 2026 au 1\textsuperscript{er} juillet 2027), alignée sur le calendrier ;
  les questions 4 et 5 tiennent compte du test de positionnement du 29 septembre 2026 (voir le guide et l'annexe A).
\else
  Chaque fiche comprend cinq questions à traiter en sept minutes, sans calculatrice sauf mention.
  Les questions 4 et 5 reviennent sur des notions déjà rencontrées : elles te permettent de te remettre en mémoire ce que tu as déjà appris.
  Indique toujours ta méthode, même pour un calcul mental.
\fi
\end{mdframed}
\tableofcontents
\bigskip
\foreach\i in {01,02,03,04,05,06,07,08,09,10,11,12,13,14}{\input{sections/ch\i}}
\ifprof
  \clearpage
  % Guide d'utilisation — version professeur uniquement
\section{Guide d'utilisation (version professeur)}
\subsection*{Principe}
Chaque séance de 2G ouvre sur une fiche de cinq questions (5 à 8 minutes), alignée sur la séance du calendrier du jour :
\begin{enumerate}[label=\textbf{\arabic*.},leftmargin=1.6em]
  \item \textbf{Échauffement} : un prérequis de cycle~4 directement utile à la séance.
  \item \textbf{Cœur} : une question sur la notion du jour (ou de la veille).
  \item \textbf{Cœur +} : une seconde question plus exigeante (QCM, vrai/faux justifié, lecture de code, méthode).
  \item \textbf{Retour} : une question hors chapitre, issue d'un attendu « à consolider » ou « à entretenir » (espacement des révisions).
  \item \textbf{Fil rouge} : une question ciblée sur un attendu en « renforcement prioritaire » du test de positionnement du 29/09/2026.
\end{enumerate}
En version professeur, chaque question est précédée de son étiquette de diagnostic (par exemple \tagdiag{EQ}) ; la réponse attendue figure en rouge, suivie d'une remarque ou d'un piège fréquent.

\subsection*{Déroulement suggéré (7 min)}
\begin{enumerate}[label=\textbf{\arabic*.},leftmargin=1.6em]
  \item 3 min : réponses individuelles à l'ardoise ou sur le cahier de rituels, sans calculatrice sauf mention.
  \item 2 min : mise en commun ; faire verbaliser une réponse \emph{fausse} et une réponse \emph{juste} (« comment as-tu fait\,? »).
  \item 2 min : correction ciblée des questions 4--5 (celles du diagnostic) avec la phrase-méthode associée.
\end{enumerate}

\subsection*{Variante FLE/FLS}
\begin{itemize}[leftmargin=1.2em,itemsep=2pt]
  \item \textbf{Consigne} : projeter la question avec le verbe d'action surligné (\emph{calculer, résoudre, montrer, comparer}) et un exemple de réponse modèle la première semaine.
  \item \textbf{Phrases-cadres} : « Je calcule\ldots{} donc\ldots{} », « Je compare\ldots{} car\ldots{} », « Le contre-exemple est\ldots{} ». Les afficher en bandeau permanent.
  \item \textbf{Glossaire visuel} : fiche mots-clés (déterminant, colinéaire, médiane, taux d'évolution, fréquence conditionnelle) avec schéma et traduction dans la langue d'origine de l'élève si utile.
  \item \textbf{Binômes} : 30 secondes de reformulation entre pairs avant de répondre ; l'élève allophone peut d'abord répondre par un geste ou un schéma, puis formaliser à l'oral.
  \item \textbf{Écrits} : accepter les notations symboliques avant la phrase rédigée, puis exiger la phrase de conclusion une fois sur deux.
\end{itemize}

\subsection*{Grille d'analyse des erreurs}
{\footnotesize\renewcommand{\arraystretch}{1.2}
\begin{tabular}{@{}p{2.2cm} p{4.4cm} p{4.8cm} p{4.6cm}@{}}\toprule
Étiquettes & Erreur observée & Indice dans la réponse & Geste de remédiation \\ \midrule
EQ, CAL & Erreur de signe en transposant ou en distribuant & $5x-7=2x+8\Rightarrow3x=1$ & Écrire chaque étape ; vérifier en remplaçant. \\
DEV, LIT & Distributivité oubliée ($2(x+3)=2x+3$) & Résultat non réduit ou faux pour $x=1$ & Test numérique ; aire d'un rectangle. \\
PUI, SUB & Confusion $(-a)^2$ et $-a^2$ & Réponse négative pour un carré & Parenthèses systématiques à la substitution. \\
GEO, REP & Distance sans racine ; $x_A-x_B$ au lieu de $x_B-x_A$ & $AB=25$ ou signe incohérent & Vérifier l'ordre de grandeur sur un schéma. \\
PRO & Addition de taux successifs ; confusion proportion/évolution & $+20\,\%$ puis $-10\,\%=+10\,\%$ & Passer par les coefficients multiplicateurs. \\
PRB, LEC & Confusion $P_A(B)$ et $P_B(A)$ & Mauvais total de référence & Entourer le total du sous-ensemble « sachant que ». \\
REG, IMG & Confusion image et antécédent ; lecture inversée & Réponse sur le mauvais axe & Colorier : abscisse = départ, ordonnée = arrivée. \\
DEM, LIT & Un exemple pris pour une preuve & « Ça marche pour $n=3$ » & Distinguer test, contre-exemple et démonstration. \\
ALG & Exécution de code sans suivre les variables & Valeur finale incorrecte & Tableau de suivi (variable / valeur / étape). \\ \bottomrule
\end{tabular}}

\subsection*{Barème}
Les rituels sont formatifs : aucun barème n'est joint. Un barème détaillé par compétence (Chercher, Modéliser, Représenter, Raisonner, Calculer, Communiquer) peut être ajouté si l'une de ces fiches sert de support d'évaluation.

\subsection*{Sources et points de vigilance}
\begin{itemize}[leftmargin=1.2em,itemsep=2pt]
  \item Programme de mathématiques de seconde générale et technologique, arrêté du 26-2-2026 (NOR MENE2602914A), BO n°14 du 2 avril 2026 : rubrique « Automatismes » et parties « Nombres et calculs », « Géométrie », « Fonctions », « Statistiques et probabilités », « Vocabulaire ensembliste et logique ».
  \item Calendrier : consulté le 06/10/2026 ; l'ordre des notions suit le calendrier et non les dates du document de progression, qui diffèrent (ex.\ fonctions : 3--15 déc.\ 2026 au calendrier).
  \item Exercices : \textbf{créés pour ce document} avec votre autorisation ; toutes les valeurs numériques et algébriques ont été recalculées par script (voir \texttt{tools/}), les raisonnements ont été relus. Les données du tableau « précipitations » (fiche~35) sont fictives.
  \item Convention des quartiles : plus petite valeur dont au moins $25\,\%$ (resp.\ $75\,\%$) des données lui sont inférieures ou égales (rang $\lceil n/4\rceil$, $\lceil3n/4\rceil$) --- \textbf{[VERIFY]} coordonner avec les conventions du manuel utilisé.
  \item Les questions sur les solides, patrons et perspective (étiquettes SOL, PER) relèvent du cycle~4 : le programme de Seconde ne comporte pas de chapitre dédié ; elles servent le diagnostic.
  \item Fiche~77 (question~3) : notion de taux d'accroissement en anticipation de la Première, hors chapitre de Seconde.
  \item Liste des automatismes évaluables à l'épreuve anticipée (BO n°24 du 12 juin 2025, session 2027) : \textbf{[VERIFY --- check Eduscol]} pour la session 2028.
\end{itemize}

  \clearpage
  % Généré par tools/build.py
\section{Annexe A — Diagnostic du 29 septembre 2026 et couverture par les fiches}
Données : test de positionnement 2G du 29/09/2026 (moyenne de classe en \%, nombre d'élèves sous $50\,\%$ et au-dessus de $80\,\%$). La colonne « Fiches » indique le nombre de fiches (sur $88$) où l'étiquette de diagnostic correspondante apparaît au moins une fois ; plusieurs attendus partagent une même étiquette.\par\smallskip
{\footnotesize\renewcommand{\arraystretch}{1.15}\begin{tabular}{@{}l l p{7.6cm} r r r l r@{}}\toprule
N° & Étiq. & Attendu & Moy. & $<50\%$ & $\geqslant80\%$ & Priorité & Fiches \\ \midrule
01 & CON & Comprendre un algorithme — programme de construction & 60,7 & 11 & 17 & À entretenir & 6 \\
02 & ALG & Comprendre un algorithme — programme de calcul & 42,9 & 16 & 12 & \textbf{Renf.\ prioritaire} & 17 \\
03 & NUM & Diverses représentations d'un même nombre (décimale, fractionnaire, notation scientifique, droite graduée) & 74,5 & 3 & 14 & À entretenir & 13 \\
04 & CAL & Calculer avec relatifs, fractions, décimaux & 65,4 & 6 & 8 & À entretenir & 21 \\
05 & PUI & Calculs simples avec puissances, notation scientifique & 67,1 & 7 & 15 & À entretenir & 6 \\
06 & SUB & Substituer dans une expression algébrique & 57,1 & 4 & 2 & À consolider & 8 \\
07 & DEV & Développer, factoriser, réduire (cas très simples) & 74,1 & 2 & 11 & À entretenir & 8 \\
08 & EQ & Résoudre algébriquement une équation du premier degré & 38,0 & 21 & 2 & \textbf{Renf.\ prioritaire} & 29 \\
09 & LIT & Calcul littéral~: démontrer, valider/réfuter une conjecture, traduire par une expression & 34,8 & 15 & 0 & \textbf{Renf.\ prioritaire} & 13 \\
10 & PRO & Problèmes de proportionnalité (pourcentages, échelles, agrandissement-réduction) & 52,0 & 13 & 1 & À consolider & 18 \\
11 & QP & Calculer une quatrième proportionnelle & 64,3 & 10 & 18 & À entretenir & 5 \\
12 & GEO & Figures, configurations, transformations~: déterminer des grandeurs & 35,5 & 24 & 0 & \textbf{Renf.\ prioritaire} & 28 \\
13 & STA & Effectifs, fréquences\,; moyenne, médiane, étendue & 51,8 & 8 & 9 & À consolider & 15 \\
14 & LEC & Lire et interpréter données brutes, tableaux, diagrammes & 57,1 & 12 & 6 & À consolider & 14 \\
15 & IMG & Image ou antécédent à partir d'une représentation & 76,2 & 6 & 14 & À entretenir & 12 \\
16 & MOD & Problèmes modélisés par des fonctions & 78,6 & 6 & 22 & À entretenir & 13 \\
17 & DEM & Raisonner et s'initier à la démonstration (figures, configurations) & 34,2 & 23 & 1 & \textbf{Renf.\ prioritaire} & 17 \\
18 & REP & Repérage (droite graduée, plan, pavé, sphère) & 32,1 & 19 & 9 & \textbf{Renf.\ prioritaire} & 15 \\
19 & SOL & Reconnaître des solides & 34,5 & 20 & 1 & \textbf{Renf.\ prioritaire} & 10 \\
20 & PNP & Reconnaître proportionnalité / non-proportionnalité & 61,9 & 10 & 8 & À entretenir & 5 \\
21 & CMP & Calculs et comparaisons pour traiter des problèmes & 52,7 & 6 & 2 & À consolider & 13 \\
22 & DIV & Problèmes de divisibilité (engrenages, conjonction de phénomènes) & 42,9 & 10 & 6 & \textbf{Renf.\ prioritaire} & 11 \\
23 & PER & Représentations de solides (perspective, face/dessus, sections, patrons) & 67,9 & 9 & 19 & À entretenir & 5 \\
24 & PRB & Probabilités en situation d'équiprobabilité & 57,1 & 12 & 16 & À consolider & 17 \\
25 & REG & Passer d'un mode de représentation d'une fonction à un autre & 25,0 & 21 & 7 & \textbf{Renf.\ prioritaire} & 21 \\
\bottomrule\end{tabular}}\par\bigskip
\textbf{Fréquence des questions « Retour » et « Fil rouge » par groupe de priorité} (nombre de fiches contenant au moins une question de ce groupe en positions 4–5) :\par\smallskip
\begin{tabular}{@{}l r@{}}\toprule Groupe & Fiches sur 88 \\ \midrule
Renforcement prioritaire & 88 \\
À consolider & 36 \\
À entretenir & 52 \\
\bottomrule\end{tabular}\par
\fi
\end{document}
'` ou en remplaçant \proffalse par \proftrue
\ifdefined\PROFVERSION\proftrue\fi

% --- Compteurs (automatiques) ----------------------------------------
\newcounter{fiche}
\newcounter{qfiche}[fiche]

% --- Étiquette de diagnostic (version professeur) --------------------
\newcommand{\tagdiag}[1]{\textcolor{acier}{\footnotesize\textsf{[#1]}}}

% --- Conteneur d'une fiche ------------------------------------------
\mdfdefinestyle{fichestyle}{%
  linewidth=2pt,topline=false,bottomline=false,rightline=false,
  leftline=true,linecolor=marine,innerleftmargin=9pt,innerrightmargin=2pt,
  innertopmargin=4pt,innerbottommargin=5pt,skipabove=7pt,skipbelow=5pt,
  backgroundcolor=white}
\mdfsetup{style=fichestyle}
\makeatletter
\newcommand{\fichehead}[2]{\gdef\fh@date{#1}\gdef\fh@titre{#2}}
\newenvironment{fiche}{%
  \stepcounter{fiche}\Needspace{9\baselineskip}\begin{mdframed}[style=fichestyle]%
  \noindent{\color{marine}\textbf{\textsf{Fiche~\thefiche}}}\hspace{.6em}%
  {\color{ardoise}\small\textsf{\fh@date}}\hspace{.6em}{\itshape\small\fh@titre}\par\vspace{3pt}%
  \begin{enumerate}[label=\textbf{\arabic*.},leftmargin=1.5em,itemsep=3pt,parsep=0pt]}{%
  \end{enumerate}\end{mdframed}}
\makeatother
\newcommand{\role}[1]{{\color{acier}\scriptsize\textsf{#1}}}
\newcommand{\qitem}[5]{%
  \item \role{#1}\ \ifprof\tagdiag{#2}\fi\ #3%
  \ifprof
    \par\smallskip\noindent{\color{rep}\small\textbf{Rép.}\ #4}%
    \if\relax\detokenize{#5}\relax\else
      \par\noindent{\color{ardoise}\footnotesize\textit{#5}}\fi
  \fi}
\newcommand{\choix}[4]{\par\smallskip\noindent\begin{tabular}{@{}*{4}{p{.22\linewidth}}@{}}
  \textbf{A.}~#1 & \textbf{B.}~#2 & \textbf{C.}~#3 & \textbf{D.}~#4\end{tabular}}

% --- Titres, en-têtes ------------------------------------------------
\titleformat{\section}{\Large\sffamily\bfseries\color{marine}}{}{0pt}{}[\vspace{-2pt}{\color{acier}\titlerule}]
\titlespacing*{\section}{0pt}{12pt}{6pt}
\pagestyle{fancy}\fancyhf{}
\renewcommand{\headrulewidth}{.3pt}
\fancyhead[L]{\small\sffamily Rituels de début de cours — 2G}
\fancyhead[R]{\small\sffamily\ifprof Version professeur\else Version élève\fi}
\fancyfoot[C]{\small\thepage}
\setlength{\parindent}{0pt}\raggedright\hyphenpenalty=10000\exhyphenpenalty=10000
\usepackage[hidelinks]{hyperref}  % toujours en dernier (ou presque)
